
\documentclass[11pt]{article}

\usepackage[a4paper,margin=1in]{geometry}
\usepackage{amsmath,amssymb,amsthm,mathtools}
\usepackage{booktabs}
\usepackage{array}
\usepackage{enumitem}
\usepackage{hyperref}
\usepackage{microtype}
\usepackage{xcolor}
\usepackage{listings}

\hypersetup{
    hidelinks,
    pdftitle={From Distribution-Free CDF Bands to Polyhedral Confidence Regions for Location--Scale Families}
}

\definecolor{codecomment}{RGB}{40,110,65}
\definecolor{codekeyword}{RGB}{35,70,150}
\definecolor{codebackground}{RGB}{247,248,250}

\lstdefinestyle{Rcode}{
    language=R,
    basicstyle=\ttfamily\scriptsize,
    keywordstyle=\color{codekeyword}\bfseries,
    commentstyle=\color{codecomment}\itshape,
    stringstyle=\color{purple!70!black},
    backgroundcolor=\color{codebackground},
    frame=single,
    rulecolor=\color{black!20},
    numbers=left,
    numberstyle=\tiny\color{black!50},
    stepnumber=1,
    numbersep=8pt,
    breaklines=true,
    breakatwhitespace=false,
    showstringspaces=false,
    keepspaces=true,
    columns=fullflexible,
    tabsize=2,
    xleftmargin=1.5em,
    framexleftmargin=1.2em,
    aboveskip=0.8\baselineskip,
    belowskip=0.8\baselineskip
}

\newtheorem{proposition}{Proposition}
\newtheorem{remark}{Remark}
\newtheorem{lemma}{Lemma}

\newcommand{\R}{\mathbb{R}}
\newcommand{\Pp}{\mathbb{P}}
\newcommand{\Ee}{\mathbb{E}}
\newcommand{\cC}{\mathcal{C}}
\newcommand{\cB}{\mathcal{B}}
\newcommand{\ind}{\mathbf{1}}

%\title{
%From Distribution-Free CDF Bands to Polyhedral Confidence Regions\\
%for Location--Scale Families
%}

%\author{[Author Name]}
%\date{\today}

\begin{document}

%\maketitle

% ----------------------------------------------------------------------
% FRONT MATTER
% ----------------------------------------------------------------------

\begin{titlepage}

\centering

\vspace*{1.5cm}

{\LARGE\bfseries
From Distribution-Free CDF Bands to Polyhedral Confidence Regions\\[0.3em]
for Location--Scale Families
\par}

\vspace{1.0cm}

{\large
A documented AI-generated research artifact
\par}

\vspace{1.5cm}

{\large\bfseries Roman Guchenko\par}
\vspace{0.2cm}
{\normalsize
Human initiator, curator, and verifier
\par}

\vspace{1.2cm}

{\normalsize
Substantive scientific development and initial manuscript generation
\par}
\vspace{0.2cm}
{\large\bfseries
ChatGPT (OpenAI, GPT-5.6 Sol)
\par}

\vspace{1.5cm}

\begin{minipage}{0.90\textwidth}
\small
\textbf{Research provenance.}
This work arose from an extended research dialogue between the named
human participant and ChatGPT.  The human participant posed the initial
research question, selected among directions proposed during the
investigation, executed computational experiments, returned their
results to the model, later personally reviewed and verified the
mathematical arguments and R implementation, curated the resulting
research record, and prepared the work for public dissemination.

ChatGPT generated most of the subsequent substantive scientific
development, including literature exploration, formulation and
generalization of the statistical framework, mathematical derivations,
algorithm design, R code drafts, simulation design, interpretation of
results, and the initial manuscript text.  The named human therefore
does not claim conventional authorship of those AI-generated
contributions.  His role in this publication is stated explicitly as
human initiator, curator, and verifier.

The research-development transcript, source code, computational
notebook, verification materials, and manuscript source are preserved
with this paper so that the provenance of the work can be examined
directly.
\end{minipage}

\vfill

{\small
Research developed August 2026\\
Public version: 20 September 2026\\
Research record: \url{https://github.com/RomanGuchenko/cdf-band-location-scale-regions}
\par}

\end{titlepage}

\begin{abstract}
This paper studies confidence regions for the location and scale parameters of a
continuous location--scale family obtained by inverting simultaneous,
distribution-free confidence bands for the cumulative distribution
function.  If the band implies simultaneous order-statistic constraints
\[
    \ell_i \leq F(X_{(i)}) \leq u_i,
    \qquad i=1,\ldots,n,
\]
then for a location--scale family
\[
    F_{\mu,\sigma}(x)
    =
    G\!\left(\frac{x-\mu}{\sigma}\right),
    \qquad \sigma>0,
\]
the induced confidence region is the intersection of the affine
constraints
\[
    X_{(i)}-\sigma G^{-1}(u_i)
    \leq \mu
    \leq
    X_{(i)}-\sigma G^{-1}(\ell_i).
\]
Consequently, the region is convex and polyhedral in the
$(\mu,\sigma)$ parameterization.  Exact representations are derived for its
projection onto the scale and location coordinates, its piecewise-linear
boundary, and its area.  When the band limits are monotone in the order
statistic index, the complete boundary can be constructed by a
monotone-slope envelope algorithm in linear time after sorting the
sample.

The framework permits different distribution-free CDF bands to be
compared according to their induced parametric uncertainty rather than
according to CDF-band width alone.  The study investigates Kolmogorov--Smirnov,
Dvoretzky--Kiefer--Wolfowitz, Berk--Jones/Owen, and
Dümbgen--Wellner bands.  Preliminary Monte Carlo experiments reveal a
strong target-dependent pattern.  For the normal location--scale family,
the Kolmogorov--Smirnov band gives the shortest marginal interval for
location, Berk--Jones gives the shortest interval for scale, and the
Dümbgen--Wellner band gives the smallest expected joint area across a
range of sample sizes.  Experiments with Laplace and Student-$t$
families show that the magnitude, and in some cases the ranking, of these
efficiency differences depends on the base location--scale family.
These observations motivate the broader problem of designing
distribution-free CDF bands specifically to optimize induced parametric
confidence regions.
\end{abstract}


\section*{Research provenance, attribution, and verification}

\subsection*{Origin of the investigation}

This paper has an unusual research provenance.  The project began with
a question posed by the human participant in a ChatGPT conversation:
whether confidence information for a normal cumulative distribution
function, viewed relative to an empirical CDF, could be inverted to
obtain confidence information for the normal location and scale
parameters.

ChatGPT initially recognized the question as one of test inversion and
related it to distribution-free goodness-of-fit procedures.  As the
dialogue continued, the model observed that inversion of suitable
order-statistic constraints produces affine inequalities in location
and scale.  It then generalized this observation from the normal model
and the Kolmogorov--Smirnov procedure to arbitrary continuous
location--scale families and, subsequently, to arbitrary simultaneous
distribution-free order-statistic CDF bands.

The human participant did not prescribe this mathematical development
in advance.  His role during the research dialogue was principally to
pose questions, choose among directions suggested by the model, request
deeper investigation of particular possibilities, execute proposed
computations, return numerical results, and decide which developments
were worth pursuing.

\subsection*{Development of the research}

The scientific direction changed repeatedly during the conversation.
Early candidate contributions were subjected to literature searches
performed by ChatGPT.  In particular, the model initially identified
convexity of the Kolmogorov--Smirnov-inverted location--scale region as
potentially interesting, but subsequent literature exploration showed
that this property and closely related constructions were already
known.  The research direction was therefore narrowed rather than
retaining the earlier novelty claim.

A further literature search identified recent work on inversion of
distribution-free CDF bands through parametric quantile functions.
This again changed the scope of the project.  The investigation moved
toward the special geometric structure available for two-parameter
location--scale families: affine constraints, exact projections,
piecewise-linear boundaries, exact region area, location--scale
equivariance, and efficient envelope algorithms.  The dialogue then
developed the comparison of different CDF bands according to the
parametric confidence regions that they induce.

ChatGPT subsequently generated the principal mathematical derivations,
the envelope-based computational construction, drafts of the R
implementation, the Monte Carlo design, interpretations of the
simulation results, and the initial version of the present manuscript.
The research record also contains corrections and changes of direction
made during this process.  It is therefore preserved as a chronological
record rather than being rewritten as a retrospective account in which
the final argument appears to have been known from the beginning.

\subsection*{Human verification}

After the AI-generated development of the work, the human curator
personally reviewed and verified the mathematical arguments and
implementation before deciding to disseminate the research.

The accompanying computational materials additionally contain a
separate verification suite.  Among other checks, it compares the
linear-time construction with an independent quadratic pairwise
formula, compares the computed envelopes with direct evaluations,
checks exact area calculations against numerical integration, tests
affine equivariance, exercises empty and unbounded cases, verifies that
band inversion reproduces the originating test event, checks seeded
calibration reproducibility, and includes an empirical diagnostic for
the expected computational scaling.

These automated checks are part of the reproducibility record.  They
are distinct from the human curator's responsibility for examining the
mathematics and implementation before publication.

\subsection*{Attribution}

The division of contributions in this project does not fit ordinary
single-author terminology.  The substantive mathematical and
computational research was generated predominantly by ChatGPT, while
the human participant initiated the investigation, made choices about
its direction, executed experiments, verified the resulting work,
preserved the record, and arranged its dissemination.

For that reason, Roman Guchenko is identified on the title page as
\emph{human initiator, curator, and verifier}, rather than being
presented without qualification as the scientific author of the
results.  ChatGPT is explicitly credited with the substantive
scientific development and initial manuscript generation.

This attribution is intended to describe the provenance of the work as
accurately as possible.  It does not imply that the research was
produced autonomously without human interaction: the complete record
shows an iterative dialogue in which the human participant repeatedly
selected questions and supplied computational feedback.  Conversely,
the human role in that dialogue should not be interpreted as a claim
that the principal mathematical constructions, derivations, algorithms,
or initial manuscript were originated by the human participant.

\subsection*{Public research record}

The intended public research record consists of the following
components:

\begin{itemize}
    \item the present technical manuscript and its \LaTeX{} source;
    \item a reader-oriented rendering of the chronological
          human--AI research transcript;
    \item an archival version of the retrieved transcript;
    \item the R reference implementation;
    \item the separate R verification suite;
    \item the computational notebook and its rendered form; and
    \item the information required to reproduce the reported
          simulations, including seeds and calibration settings.
\end{itemize}

The transcript is provided not merely as supplementary commentary.
Because the unusual provenance of the scientific work is itself an
important claim of this publication, the dialogue constitutes primary
evidence for how the research developed.

At the time this version was prepared, the retrieved conversation
record contained one explicitly identified truncation caused by a
source-interface length limit.  If the corresponding passage can be
recovered from the original conversation export, it should replace the
truncated passage in the archival record.  Otherwise, that limitation
will remain explicitly documented rather than being silently repaired.


\section{Introduction}

A simultaneous confidence band for a cumulative distribution function
(CDF) is usually viewed as a nonparametric inferential object.  Given an
empirical distribution function $F_n$, a band provides a set of
distribution functions that are compatible with the observed sample at
a prescribed confidence level.  If a parametric family
$\{F_\theta:\theta\in\Theta\}$ is of interest, however, one may instead
ask which members of this parametric family are compatible with the
same band.  This leads naturally to the inverted set
\[
    \cC(X)
    =
    \left\{
        \theta\in\Theta:
        F_\theta \text{ lies in the simultaneous CDF band}
    \right\}.
\]
The resulting set is a confidence region for the parameter.

Inference obtained by inversion of goodness-of-fit procedures is
classical.  Easterling \cite{Easterling1976} explicitly proposed
parameter estimation through inversion of goodness-of-fit tests, while
Littell and Rao \cite{LittellRao1978} developed confidence regions for
location and scale parameters based on the Kolmogorov--Smirnov (KS)
statistic.  Salvia \cite{Salvia1980} subsequently established important
geometric properties of KS consonance sets, including convexity and
finiteness for location--scale models.  Thus, neither goodness-of-fit
test inversion nor convexity of the KS location--scale region is new.

At the same time, the literature on distribution-free CDF bands has
developed substantially beyond KS.  Berk and Jones
\cite{BerkJones1979} introduced likelihood-ratio-type goodness-of-fit
statistics with strong tail sensitivity.  Owen \cite{Owen1995} inverted
the Berk--Jones construction to obtain nonparametric likelihood
confidence bands, which are narrower in the tails and wider in the
center than KS bands.  Frey \cite{Frey2008} studied distribution-free
bands chosen directly to optimize narrowness criteria rather than test
power.  More recently, Dümbgen and Wellner \cite{DumbgenWellner2023}
introduced a multiscale refinement of the Berk--Jones/Owen framework
that preserves strong tail behavior while improving performance in the
central region.

There has also been renewed interest in direct inversion of CDF bands
through parametric quantile functions.  Chattopadhyay, Sarkar, and
Kuchibhotla \cite{Chattopadhyay2026} develop a general framework for
quantile-parametrized families, motivated particularly by models in
which the quantile function is tractable while the likelihood or CDF is
not.

The present work focuses on a special case that has unusually strong
geometric structure: two-parameter location--scale families.  The
quantile function is affine in the parameters,
\[
    Q_{\mu,\sigma}(p)
    =
    \mu+\sigma G^{-1}(p),
\]
and therefore inversion of any simultaneous order-statistic CDF band
produces linear inequalities in $(\mu,\sigma)$.  This observation allows
the induced confidence set to be treated as a problem in convex
geometry.

The objectives of this paper are fourfold.

First, the paper formulates CDF-band inversion for a generic continuous
location--scale family and shows that the resulting confidence region is
polyhedral in $(\mu,\sigma)$.

Second, exact computational representations are derived for the joint
region, its marginal confidence intervals for $\mu$ and $\sigma$, and
its area.  These calculations reduce to piecewise-linear envelope
operations.

Third, this geometry is used to compare several classical and modern
distribution-free confidence bands according to three parametric
criteria:
\[
    \operatorname{Area}(\cC),\qquad
    |\cC_\mu|,\qquad
    |\cC_\sigma|.
\]

Fourth, the study investigates how these comparisons change with the base
location--scale family.  The preliminary results suggest that
CDF-band quality is inherently target dependent: a band that is highly
efficient for scale need not be efficient for location, and a band that
minimizes joint area need not minimize either marginal length.


\section{General band inversion}

\subsection{Order-statistic representation}

Let
\[
    X_1,\ldots,X_n
\]
be independent observations from a continuous distribution $F$, and let
\[
    X_{(1)}\leq\cdots\leq X_{(n)}
\]
denote the order statistics.

Suppose that a simultaneous distribution-free CDF band implies
deterministic bounds
\begin{equation}
    \ell_i
    \leq
    F(X_{(i)})
    \leq
    u_i,
    \qquad i=1,\ldots,n,
    \label{eq:orderband}
\end{equation}
such that
\begin{equation}
    \Pp_F
    \left(
        \ell_i
        \leq
        F(X_{(i)})
        \leq
        u_i
        \text{ for every }i
    \right)
    \geq
    1-\alpha
    \label{eq:coverage}
\end{equation}
for every continuous $F$.

The distribution-free nature of \eqref{eq:coverage} follows from the
probability integral transform.  If
\[
    U_i=F(X_i),
\]
then $U_i$ are independent $\operatorname{Unif}(0,1)$ random variables
and
\[
    F(X_{(i)})=U_{(i)}.
\]
Thus calibration of the band depends only on uniform order statistics
and not on the eventual parametric family.

A small indexing detail is important.  If a CDF band is originally
specified on the intervals between empirical jumps, then at
$X_{(i)}$ continuity of $F$ combines information from the empirical
CDF immediately before and immediately after the jump.  In particular,
for constant-width bands the lower endpoint involves $i/n$ whereas the
upper endpoint involves $(i-1)/n$.  This distinction is essential for
obtaining the intended finite-sample coverage.


\subsection{Location--scale families}

Consider the location--scale family
\begin{equation}
    F_{\mu,\sigma}(x)
    =
    G\!\left(
        \frac{x-\mu}{\sigma}
    \right),
    \qquad
    \mu\in\R,\quad
    \sigma>0,
    \label{eq:lsfamily}
\end{equation}
where $G$ is a known continuous strictly increasing CDF.

Let
\[
    g(p)=G^{-1}(p),
\]
and define
\begin{equation}
    a_i=g(\ell_i),
    \qquad
    b_i=g(u_i).
    \label{eq:ab}
\end{equation}
Values $g(0)=-\infty$ and $g(1)=+\infty$ are interpreted in the usual
extended-real sense.

For a candidate $(\mu,\sigma)$, condition \eqref{eq:orderband} becomes
\[
    \ell_i
    \leq
    G\!\left(
        \frac{X_{(i)}-\mu}{\sigma}
    \right)
    \leq
    u_i.
\]
Since $G$ is increasing,
\[
    a_i
    \leq
    \frac{X_{(i)}-\mu}{\sigma}
    \leq
    b_i,
\]
and therefore
\begin{equation}
    X_{(i)}-\sigma b_i
    \leq
    \mu
    \leq
    X_{(i)}-\sigma a_i.
    \label{eq:linearconstraints}
\end{equation}

This gives the fundamental geometric representation.

\begin{proposition}[Band-induced location--scale region]
\label{prop:polyhedral}
Let \eqref{eq:coverage} hold for a simultaneous order-statistic band.
For the location--scale family \eqref{eq:lsfamily}, define
\begin{equation}
    \cC
    =
    \bigcap_{i=1}^n
    \left\{
        (\mu,\sigma):
        X_{(i)}-\sigma b_i
        \leq
        \mu
        \leq
        X_{(i)}-\sigma a_i
    \right\}
    \cap
    \{\sigma>0\}.
    \label{eq:C}
\end{equation}
Then $\cC$ is a convex polyhedral set in the $(\mu,\sigma)$
parameterization, and
\[
    \Pp_{\mu_0,\sigma_0}
    \left(
        (\mu_0,\sigma_0)\in\cC
    \right)
    \geq
    1-\alpha.
\]
If the order-statistic band is calibrated with exact coverage
$1-\alpha$, the induced region has the same coverage.
\end{proposition}

\begin{proof}
Each inequality in \eqref{eq:linearconstraints} defines a half-plane in
$(\mu,\sigma)$, and $\{\sigma>0\}$ is convex.  Their intersection is
therefore convex and polyhedral.

For the true parameter $(\mu_0,\sigma_0)$,
membership in $\cC$ is equivalent to
\[
    \ell_i
    \leq
    F_{\mu_0,\sigma_0}(X_{(i)})
    \leq
    u_i
\]
for every $i$.  The coverage statement therefore follows directly from
\eqref{eq:coverage}.
\end{proof}

\begin{remark}
Convexity here is naturally expressed in $(\mu,\sigma)$ coordinates.
The transformation from scale to variance,
$v=\sigma^2$, is nonlinear and does not in general preserve convexity.
\end{remark}


\section{Exact geometry of the confidence region}

Write
\[
    x_i=X_{(i)}
\]
for convenience.

\subsection{Lower and upper envelopes}

For any fixed scale $s>0$, define
\begin{align}
    L(s)
    &=
    \max_{1\leq i\leq n}
    \left\{
        x_i-b_i s
    \right\},
    \label{eq:L}
    \\
    U(s)
    &=
    \min_{1\leq i\leq n}
    \left\{
        x_i-a_i s
    \right\}.
    \label{eq:U}
\end{align}
Constraints involving infinite $a_i$ or $b_i$ are omitted whenever they
are nonbinding.

The joint region can be written exactly as
\begin{equation}
    \cC
    =
    \left\{
        (\mu,s):
        s>0,\quad
        L(s)\leq\mu\leq U(s)
    \right\}.
    \label{eq:enveloperegion}
\end{equation}

Because $L$ is a maximum of affine functions, it is convex and
piecewise linear.  Because $U$ is a minimum of affine functions, it is
concave and piecewise linear.  Consequently
\[
    H(s)=U(s)-L(s)
\]
is concave and piecewise linear.

The feasible scale values satisfy
\[
    H(s)\geq0.
\]
A superlevel set of a concave function on the real line is an interval.
Thus, whenever the joint confidence region is nonempty and bounded in
scale,
\begin{equation}
    \cC_\sigma
    =
    [\sigma_L,\sigma_U].
\end{equation}


\subsection{Pairwise formula for the scale interval}

For a fixed $s$, feasibility requires every lower constraint to lie
below every upper constraint:
\[
    x_i-sb_i
    \leq
    x_j-sa_j
    \qquad
    \text{for all }i,j.
\]
Equivalently,
\begin{equation}
    s(a_j-b_i)
    \leq
    x_j-x_i.
    \label{eq:pairwise}
\end{equation}

Let
\[
    d_{ij}=a_j-b_i,
    \qquad
    \Delta_{ij}=x_j-x_i.
\]
If $d_{ij}>0$, \eqref{eq:pairwise} gives
\[
    s
    \leq
    \frac{\Delta_{ij}}{d_{ij}}.
\]
If $d_{ij}<0$, it gives
\[
    s
    \geq
    \frac{\Delta_{ij}}{d_{ij}}.
\]

Hence an exact representation is
\begin{equation}
    \sigma_L
    =
    \max
    \left\{
        0,\,
        \max_{d_{ij}<0}
        \frac{x_j-x_i}{a_j-b_i}
    \right\},
    \label{eq:sigmal}
\end{equation}
and
\begin{equation}
    \sigma_U
    =
    \min_{d_{ij}>0}
    \frac{x_j-x_i}{a_j-b_i},
    \label{eq:sigmau}
\end{equation}
with the usual conventions if one of the index sets is empty.

This representation shows that critical scale values are ratios of
observed order-statistic spacings to theoretical quantile spacings.


\subsection{Marginal interval for location}

The marginal confidence interval for $\mu$ is the projection of
\eqref{eq:enveloperegion} onto the location coordinate:
\begin{equation}
    \cC_\mu
    =
    [\mu_L,\mu_U],
\end{equation}
where
\begin{align}
    \mu_L
    &=
    \min_{\sigma_L\leq s\leq\sigma_U}
    L(s),
    \label{eq:mul}
    \\
    \mu_U
    &=
    \max_{\sigma_L\leq s\leq\sigma_U}
    U(s).
    \label{eq:muu}
\end{align}

Since $L$ is convex and piecewise linear, its minimum occurs at an
endpoint or a breakpoint of its upper envelope.  Similarly, since
$U$ is concave and piecewise linear, its maximum occurs at an endpoint
or breakpoint of its lower envelope.  The marginal interval can
therefore be obtained exactly from finitely many line intersections,
without numerical optimization over a dense parameter grid.


\subsection{Area of the joint region}

The area of the joint region is
\begin{equation}
    A(\cC)
    =
    \int_{\sigma_L}^{\sigma_U}
    \left[
        U(s)-L(s)
    \right]\,ds.
    \label{eq:area}
\end{equation}

Let
\[
    \sigma_L=s_0<s_1<\cdots<s_K=\sigma_U
\]
be the merged set of breakpoints of $L$ and $U$.  On each interval
$[s_k,s_{k+1}]$,
\[
    L(s)=c_k+d_ks,
    \qquad
    U(s)=e_k+f_ks.
\]
Therefore
\begin{align}
    A(\cC)
    =
    \sum_{k=0}^{K-1}
    \Bigg[
        &(e_k-c_k)(s_{k+1}-s_k)
        \nonumber\\
        &+
        \frac{f_k-d_k}{2}
        \left(
            s_{k+1}^2-s_k^2
        \right)
    \Bigg].
    \label{eq:areaexact}
\end{align}

Equivalently, because the gap $U-L$ is affine between consecutive
breakpoints, trapezoidal integration on the breakpoint grid is exact.


\subsection{Vertices}

Boundary lines have the form
\[
    \mu=x_i-b_i\sigma
\]
or
\[
    \mu=x_j-a_j\sigma.
\]
A lower and upper boundary line intersect when
\[
    x_i-b_i\sigma
    =
    x_j-a_j\sigma,
\]
giving
\begin{equation}
    \sigma_{ij}
    =
    \frac{x_j-x_i}{a_j-b_i},
    \qquad
    \mu_{ij}
    =
    x_i-b_i\sigma_{ij}.
    \label{eq:vertices}
\end{equation}
Candidate vertices are retained only if they satisfy all constraints.


\section{Efficient envelope construction}

The naive pairwise formulas involve $O(n^2)$ candidate intersections.
The order-statistic structure permits a more efficient construction for
the common case in which the confidence-band limits are monotone.

Suppose
\[
    \ell_1\leq\cdots\leq\ell_n,
    \qquad
    u_1\leq\cdots\leq u_n.
\]
Then, because $g=G^{-1}$ is increasing,
\[
    a_1\leq\cdots\leq a_n,
    \qquad
    b_1\leq\cdots\leq b_n.
\]

The lines defining $L$ are
\[
    L_i(s)=x_i-b_i s,
\]
whose slopes $-b_i$ are monotone.  The lines defining $-U$ are
\[
    -U_i(s)=-x_i+a_i s,
\]
whose slopes $a_i$ are also monotone.

\begin{proposition}[Monotone envelope algorithm]
\label{prop:envelope}
Assume the band limits are monotone in $i$.  Given the sorted sample,
the active lines of $L$ and $U$, together with all their breakpoints,
can be constructed in $O(n)$ time.
\end{proposition}

\begin{proof}
Consider generic lines
\[
    f_j(t)=c_j+m_jt
\]
arriving in strictly increasing slope order.  Suppose the current upper
envelope contains active lines
\[
    h_1,\ldots,h_r
\]
with increasing activation points
\[
    z_1<\cdots<z_r.
\]
A new line $q$ has slope larger than that of $h_r$.  Let $t^\ast$ be
their intersection.

If
\[
    t^\ast>z_r,
\]
then $h_r$ remains active until $t^\ast$ and $q$ becomes active
thereafter.

If instead
\[
    t^\ast\leq z_r,
\]
then $h_r$ can never be part of the new envelope.  Before $z_r$ it was
not part of the old envelope, while after $z_r$ the higher-slope line
$q$ is already above it.  Hence $h_r$ may be removed permanently.  The
same comparison is repeated with the new last active line.

Each line is pushed onto the envelope stack once and can be popped at
most once.  The total number of stack operations is therefore linear in
$n$.  Equal-slope lines are handled by retaining only the line with the
largest intercept.

For $L$, reverse the order of the lines if necessary so that the slopes
$-b_i$ are increasing.  For $U$, construct the upper envelope of
$-U_i(s)=-x_i+a_i s$ and negate it.
\end{proof}

After the two envelopes have been constructed, their sorted breakpoint
lists can be merged in another linear pass.  This yields
$\sigma_L,\sigma_U$, $\mu_L,\mu_U$, the polygon vertices, and the exact
area.

More explicitly, the merged pass evaluates
\[
    H(s)=U(s)-L(s)
\]
at every breakpoint.  Because $H$ is concave and affine between
successive breakpoints, its first and last zero crossings give
$\sigma_L$ and $\sigma_U$ exactly.  No pairwise intersection matrix is
formed.

The reference implementation uses preallocated envelope stacks and
advances active-line indices monotonically during evaluation.  Thus the
implementation, not only the abstract hull algorithm, has linear time
and storage after the observations and slopes are ordered.  The
pedagogical implementation assumes, as do the bands considered here,
that the probability limits are monotone in the order-statistic index.
A package implementation could check this condition and sort arbitrary
slopes as a fallback, at a cost of $O(n\log n)$.

For numerical stability, the implementation first centers the ordered
sample at an observed median and divides it by its maximum absolute
centered value.  The envelope construction and all tolerance decisions
are made in these normalized coordinates, using relative comparisons;
intervals, boundaries, lengths, and areas are then transformed back to
the original units.  This preserves the affine equivariance described
below even under very small or very large changes of measurement units.

Thus, conditional on already sorted observations, the complete
confidence-region calculation can be performed in $O(n)$ time.
Including sorting of the raw sample gives an overall complexity of
$O(n\log n)$.


\section{Location--scale equivariance}

The geometry has a useful invariance property.

\begin{proposition}[Equivariance]
\label{prop:equivariance}
Suppose
\[
    X_i=\mu_0+\sigma_0 Z_i,
    \qquad
    Z_i\sim G,
\]
with $\sigma_0>0$.  Let $\cC_X$ be the confidence region obtained from
the sample $X$, and let $\cC_Z$ be the corresponding region obtained
from the standardized sample $Z$.  Then
\begin{equation}
    \cC_X
    =
    \left\{
        (\mu_0+\sigma_0 m,\sigma_0 s):
        (m,s)\in\cC_Z
    \right\}.
    \label{eq:equivariance}
\end{equation}
Consequently,
\begin{align}
    \operatorname{Area}(\cC_X)
    &=
    \sigma_0^2
    \operatorname{Area}(\cC_Z),
    \\
    |\cC_{\mu,X}|
    &=
    \sigma_0
    |\cC_{\mu,Z}|,
    \\
    |\cC_{\sigma,X}|
    &=
    \sigma_0
    |\cC_{\sigma,Z}|.
\end{align}
\end{proposition}

\begin{proof}
Substitute
\[
    x_i=\mu_0+\sigma_0z_i,
    \quad
    \mu=\mu_0+\sigma_0m,
    \quad
    \sigma=\sigma_0s
\]
into \eqref{eq:linearconstraints}.  After cancellation and division by
$\sigma_0$,
\[
    z_i-sb_i
    \leq
    m
    \leq
    z_i-sa_i,
\]
which is precisely the standardized confidence region.
The scaling statements follow from the affine transformation and its
Jacobian.
\end{proof}

This result has an important computational implication: for a given
base family $G$, sample size $n$, confidence level, and CDF band, the
distribution of the normalized confidence-region geometry can be
studied entirely under
\[
    \mu_0=0,
    \qquad
    \sigma_0=1.
\]


\section{Distribution-free bands considered}

Four simultaneous distribution-free CDF bands are compared.  In each
case the band is calibrated on the uniform probability scale and then
converted to order-statistic bounds $(\ell_i,u_i)$.


\subsection{Kolmogorov--Smirnov}

Let
\[
    D_n
    =
    \sup_x
    |F_n(x)-F(x)|.
\]
Let $d_{n,\alpha}$ satisfy
\[
    \Pp(D_n\leq d_{n,\alpha})=1-\alpha
\]
for continuous $F$.  It may be evaluated exactly or calibrated by
simulation from independent uniform observations.

Using
\[
    D_n^+
    =
    \max_i
    \left(
        \frac{i}{n}-U_{(i)}
    \right),
\]
and
\[
    D_n^-
    =
    \max_i
    \left(
        U_{(i)}-\frac{i-1}{n}
    \right),
\]
the order-statistic bounds are
\begin{align}
    \ell_i^{\mathrm{KS}}
    &=
    \max
    \left\{
        0,\,
        \frac{i}{n}-d_{n,\alpha}
    \right\},
    \label{eq:kslower}
    \\
    u_i^{\mathrm{KS}}
    &=
    \min
    \left\{
        1,\,
        \frac{i-1}{n}+d_{n,\alpha}
    \right\}.
    \label{eq:ksupper}
\end{align}


\subsection{Dvoretzky--Kiefer--Wolfowitz}

The Dvoretzky--Kiefer--Wolfowitz inequality, with Massart's sharp
constant \cite{DKW1956,Massart1990}, gives
\[
    \Pp
    \left(
        \sup_x |F_n(x)-F(x)|>\epsilon
    \right)
    \leq
    2e^{-2n\epsilon^2}.
\]
Taking
\begin{equation}
    \epsilon_{n,\alpha}
    =
    \sqrt{
        \frac{\log(2/\alpha)}{2n}
    },
\end{equation}
gives
\begin{align}
    \ell_i^{\mathrm{DKW}}
    &=
    \max
    \left\{
        0,\,
        \frac{i}{n}-\epsilon_{n,\alpha}
    \right\},
    \\
    u_i^{\mathrm{DKW}}
    &=
    \min
    \left\{
        1,\,
        \frac{i-1}{n}+\epsilon_{n,\alpha}
    \right\}.
\end{align}

Unlike exactly calibrated KS, DKW is an inequality-based guarantee and
is generally conservative in finite samples.


\subsection{Berk--Jones/Owen}

Define the Bernoulli Kullback--Leibler divergence
\begin{equation}
    K(p,q)
    =
    p\log\frac{p}{q}
    +
    (1-p)
    \log
    \frac{1-p}{1-q},
    \label{eq:KL}
\end{equation}
with the usual endpoint conventions.

For the two-sided Berk--Jones statistic,
\begin{equation}
    T_n^{\mathrm{BJ}}
    =
    \max_{1\leq i\leq n}
    \max
    \left\{
        nK\!\left(
            \frac{i}{n},U_{(i)}
        \right),
        nK\!\left(
            \frac{i-1}{n},U_{(i)}
        \right)
    \right\}.
    \label{eq:BJstat}
\end{equation}
Let $\kappa_{n,\alpha}^{\mathrm{BJ}}$ denote its
$(1-\alpha)$ quantile.

The lower order-statistic endpoint is the lower solution of
\begin{equation}
    nK\!\left(
        \frac{i}{n},q
    \right)
    =
    \kappa_{n,\alpha}^{\mathrm{BJ}},
    \qquad
    q\leq \frac{i}{n},
\end{equation}
whereas the upper endpoint is the upper solution of
\begin{equation}
    nK\!\left(
        \frac{i-1}{n},q
    \right)
    =
    \kappa_{n,\alpha}^{\mathrm{BJ}},
    \qquad
    q\geq \frac{i-1}{n}.
\end{equation}
The scalar equations can be solved accurately by bisection or Brent's
method.

Owen \cite{Owen1995} developed finite-sample nonparametric likelihood
confidence bands by inversion of the Berk--Jones construction.


\subsection{Dümbgen--Wellner}

The $s=1$, $\nu=1$ version of Dümbgen and Wellner is used
\cite{DumbgenWellner2023}.  Define
\begin{align}
    C(t)
    &=
    \log
    \log
    \left(
        \frac{e}{4t(1-t)}
    \right),
    \\
    D(t)
    &=
    \log
    \left(
        1+C(t)^2
    \right),
    \\
    C_\nu(t)
    &=
    C(t)+\nu D(t).
\end{align}
For $p,q\in[0,1]$, let
\begin{equation}
    C_\nu(p,q)
    =
    \min_{
        \min(p,q)\leq t\leq\max(p,q)
    }
    C_\nu(t).
\end{equation}
Define
\begin{equation}
    H_{n,\nu}(p,q)
    =
    nK(p,q)-C_\nu(p,q).
\end{equation}

The finite-sample statistic can be evaluated at the empirical jumps:
\begin{equation}
    T_{n,\nu}^{\mathrm{DW}}
    =
    \max_{1\leq i\leq n}
    \max
    \left\{
        H_{n,\nu}
        \left(
            \frac{i}{n},U_{(i)}
        \right),
        H_{n,\nu}
        \left(
            \frac{i-1}{n},U_{(i)}
        \right)
    \right\}.
\end{equation}
Let $\kappa_{n,\nu,\alpha}^{\mathrm{DW}}$ be its
$(1-\alpha)$ quantile under the uniform model.

The lower endpoint $\ell_i^{\mathrm{DW}}$ is obtained by solving
\begin{equation}
    H_{n,\nu}
    \left(
        \frac{i}{n},q
    \right)
    =
    \kappa_{n,\nu,\alpha}^{\mathrm{DW}}
\end{equation}
on the lower side, and the upper endpoint $u_i^{\mathrm{DW}}$ is
obtained from
\begin{equation}
    H_{n,\nu}
    \left(
        \frac{i-1}{n},q
    \right)
    =
    \kappa_{n,\nu,\alpha}^{\mathrm{DW}}
\end{equation}
on the upper side.

Dümbgen and Wellner show that their construction retains the strong
tail performance of Berk--Jones-type bands while improving accuracy in
the central region.


\section{Simulation study}

\subsection{Design}

The preliminary simulation study compares the four bands according to
\begin{align}
    R_A
    &=
    \Ee[
        \operatorname{Area}(\cC)
    ],
    \\
    R_\mu
    &=
    \Ee[
        \mu_U-\mu_L
    ],
    \\
    R_\sigma
    &=
    \Ee[
        \sigma_U-\sigma_L
    ].
\end{align}
Empirical joint coverage is also recorded.

The implementation defines these summaries for every replication.
An empty confidence set contributes zero area and zero marginal
lengths.  A nonempty unbounded set contributes an infinite value for
each genuinely infinite size measure.  Coverage is evaluated
independently of these size calculations, directly from the event
\[
    \ell_i\leq U_{(i)}\leq u_i
    \qquad(i=1,\ldots,n),
\]
and is therefore meaningful for empty, bounded, and unbounded sets.
No missing values are discarded when Monte Carlo means are computed:
an \texttt{NA} propagates to the corresponding reported summary, while
the numbers of empty, bounded, and unbounded regions are reported
separately.

For KS, Berk--Jones, and Dümbgen--Wellner, critical values were
calibrated using $100{,}000$ samples from
$\operatorname{Unif}(0,1)$.  DKW uses its analytic critical width.
For each outer simulation experiment, $10{,}000$ samples were generated.

The Monte Carlo critical values use an exchangeable-rank rule rather
than an interpolated sample quantile.  If $T_1,\ldots,T_B$ are the
calibration statistics, set
\[
    k=\left\lceil(B+1)(1-\alpha)\right\rceil
\]
and use the $k$th order statistic as the critical value (or $+\infty$
when $k>B$).  For a continuous statistic and an independent future
value $T_{B+1}$, exchangeability gives
\[
    \Pp\{T_{B+1}\leq T_{(k)}\}=\frac{k}{B+1}\geq 1-\alpha.
\]
Thus the finite calibration sample has an explicit marginal coverage
guarantee.  Reported outer Monte Carlo standard errors, including the
binomial standard error for empirical coverage, are conditional on the
realized calibrated thresholds.

Because of Proposition~\ref{prop:equivariance}, all data were generated
under location zero and scale one.  More specifically, for a selected
base distribution $G$, the simulations generated
\[
    U_i\sim\operatorname{Unif}(0,1),
    \qquad
    X_i=G^{-1}(U_i).
\]
This implementation also provides a useful validation: the coverage
event is distribution-free and therefore should not change with $G$,
although area and marginal lengths may change substantially.

The first experiments use the normal family.  Additional experiments
use the Laplace family and the Student-$t_2$ location--scale family.
For Student-$t_2$, $\sigma$ denotes the scale parameter, not the
standard deviation; indeed, the variance of the $t_2$ distribution is
infinite.


\subsection{Normal location--scale family}

Table~\ref{tab:normal} reports preliminary results for the normal
location--scale family.

\begin{table}[htbp]
\centering
\small
\caption{
Preliminary Monte Carlo results for the normal location--scale family.
Each row is based on $10{,}000$ outer simulations.
}
\label{tab:normal}
\begin{tabular}{
    r
    l
    r
    r
    r
    r
}
\toprule
$n$
&
Band
&
Expected area
&
Expected $\mu$ length
&
Expected $\sigma$ length
&
Coverage
\\
\midrule
20
& KS
& 2.0533
& 1.2553
& 2.7715
& 0.9538
\\
20
& DKW
& 2.3012
& 1.3137
& 2.9678
& 0.9649
\\
20
& Berk--Jones
& 1.3554
& 1.3669
& 1.7189
& 0.9511
\\
20
& Dümbgen--Wellner
& \textbf{1.3511}
& 1.3260
& 1.7609
& 0.9537
\\
\addlinespace

50
& KS
& 0.6262
& \textbf{0.7517}
& 1.4084
& 0.9498
\\
50
& DKW
& 0.6721
& 0.7765
& 1.4612
& 0.9583
\\
50
& Berk--Jones
& 0.4328
& 0.8522
& \textbf{0.8776}
& 0.9540
\\
50
& Dümbgen--Wellner
& \textbf{0.4291}
& 0.8029
& 0.9171
& 0.9528
\\
\addlinespace

100
& KS
& 0.2975
& \textbf{0.5298}
& 0.9488
& 0.9515
\\
100
& DKW
& 0.3062
& 0.5371
& 0.9626
& 0.9548
\\
100
& Berk--Jones
& 0.2036
& 0.6016
& \textbf{0.5848}
& 0.9488
\\
100
& Dümbgen--Wellner
& \textbf{0.2030}
& 0.5613
& 0.6202
& 0.9502
\\
\addlinespace

200
& KS
& 0.1430
& \textbf{0.3709}
& 0.6509
& 0.9506
\\
200
& DKW
& 0.1460
& 0.3746
& 0.6573
& 0.9527
\\
200
& Berk--Jones
& 0.1004
& 0.4293
& \textbf{0.4041}
& 0.9472
\\
200
& Dümbgen--Wellner
& \textbf{0.0988}
& 0.3942
& 0.4292
& 0.9504
\\
\bottomrule
\end{tabular}
\end{table}

Three patterns are striking.

First, KS gives the smallest expected marginal interval for $\mu$ at
every sample size considered.

Second, Berk--Jones gives the smallest expected marginal interval for
$\sigma$.  Its scale interval is approximately $38\%$ shorter than the
KS interval throughout the range $n=20$ to $200$.

Third, Dümbgen--Wellner gives the smallest expected joint area at every
sample size considered.  Relative to KS, the reduction in expected area
is approximately $34\%$ for $n=20$ and remains approximately
$31\%$ for $n=50,100,200$.

Dümbgen--Wellner appears to provide a compromise between the location
and scale criteria.  Its expected $\mu$ interval is only about
$6\%$ longer than KS over most of the range considered, while its joint
area and scale interval are substantially smaller.

DKW behaves similarly to KS but is systematically more conservative,
as expected from its inequality-based calibration.


\subsection{Laplace and heavy-tailed Student families}

Table~\ref{tab:otherfamilies} reports preliminary results for the
Laplace and Student-$t_2$ families at $n=50$.

\begin{table}[htbp]
\centering
\small
\caption{
Preliminary Monte Carlo results for Laplace and Student-$t_2$
location--scale families, $n=50$.
}
\label{tab:otherfamilies}
\begin{tabular}{
    l
    l
    r
    r
    r
    r
}
\toprule
Family
&
Band
&
Expected area
&
Expected $\mu$ length
&
Expected $\sigma$ length
&
Coverage
\\
\midrule
Laplace
& KS
& 0.9403
& \textbf{0.7712}
& 1.9236
& 0.9549
\\
Laplace
& DKW
& 0.9990
& 0.7924
& 1.9865
& 0.9610
\\
Laplace
& Berk--Jones
& 0.7878
& 0.9427
& \textbf{1.2866}
& 0.9554
\\
Laplace
& Dümbgen--Wellner
& \textbf{0.7535}
& 0.8747
& 1.3332
& 0.9534
\\
\addlinespace

Student-$t_2$
& KS
& 1.0706
& \textbf{0.9335}
& 1.8749
& 0.9549
\\
Student-$t_2$
& DKW
& 1.1361
& 0.9583
& 1.9355
& 0.9610
\\
Student-$t_2$
& Berk--Jones
& 1.0823
& 1.1456
& 1.5146
& 0.9554
\\
Student-$t_2$
& Dümbgen--Wellner
& \textbf{0.9880}
& 1.0611
& \textbf{1.5060}
& 0.9534
\\
\bottomrule
\end{tabular}
\end{table}

For the Laplace family, the qualitative normal-theory ranking remains:
KS is best for location, Berk--Jones is best for scale, and
Dümbgen--Wellner is best for joint area.  However, the magnitude of the
joint-area improvement is smaller.  The expected area of the
Dümbgen--Wellner region is about $20\%$ below that of KS, compared with
approximately $31$--$34\%$ under the normal family.

The Student-$t_2$ results show a more substantial change.  KS remains
best for location and Dümbgen--Wellner remains best for joint area, but
Dümbgen--Wellner now also slightly outperforms Berk--Jones in expected
scale length.  Moreover, the Berk--Jones expected joint area is slightly
larger than the KS area.

Thus the choice of a favorable CDF band for parametric inference
depends not only on the target criterion but also on the quantile
geometry of the base family $G$.


\section{Why the criteria disagree}

The different behavior of the marginal location and scale intervals has
a simple interpretation through the quantile function
\[
    Q_{\mu,\sigma}(p)
    =
    \mu+\sigma g(p).
\]
Its parameter derivatives are
\begin{equation}
    \frac{\partial Q_{\mu,\sigma}(p)}
         {\partial\mu}
    =
    1,
    \qquad
    \frac{\partial Q_{\mu,\sigma}(p)}
         {\partial\sigma}
    =
    g(p).
    \label{eq:derivatives}
\end{equation}

Location shifts all quantiles equally.  Scale, in contrast, is
identified through the spread of theoretical quantiles, and the
magnitude of $g(p)$ often becomes large toward the tails.

This helps explain why tail-adaptive confidence bands can improve scale
inference dramatically.  KS and DKW allocate approximately constant
vertical CDF uncertainty across the probability scale.  Berk--Jones and
Dümbgen--Wellner redistribute uncertainty, providing substantially
tighter tail bounds.  Those tail constraints can have high leverage on
$\sigma$ after transformation through $G^{-1}$.

The same observation also explains why a band optimized for scale need
not be optimal for location.  Allocating more of the simultaneous-error
budget toward the center can improve location inference, whereas tighter
tail bounds can improve scale inference.

For very heavy-tailed $G$, however, $G^{-1}(p)$ grows extremely rapidly
as $p$ approaches zero or one.  Small changes in probability-scale band
endpoints may therefore be magnified strongly in parameter space.
The Student-$t_2$ experiment suggests that this interaction can alter
the ranking of the bands themselves.


\section{Relation to existing literature and proposed contribution}

The main ingredients of this framework have important precedents.

Goodness-of-fit inversion as an estimation principle is classical
\cite{Easterling1976}.  KS-based location--scale confidence regions were
developed by Littell and Rao \cite{LittellRao1978}, and Salvia
\cite{Salvia1980} proved convexity and finiteness of KS consonance sets.
Therefore, the fact that a KS-inverted location--scale region is convex
cannot be regarded as a new result.

Similarly, the CDF bands used here are existing constructions.
Berk--Jones \cite{BerkJones1979}, Owen \cite{Owen1995}, and
Dümbgen--Wellner \cite{DumbgenWellner2023} provide the principal
tail-adaptive procedures considered in this paper.  Frey
\cite{Frey2008} studies optimal distribution-free CDF bands under
band-narrowness criteria.  Frey, Marrero, and Norton
\cite{FreyMarreroNorton2009} study minimum-area confidence sets for a
normal distribution.  Most recently, Chattopadhyay, Sarkar, and
Kuchibhotla \cite{Chattopadhyay2026} develop general inversion of
distribution-free CDF bands through parametric quantile functions.

The contribution pursued here is therefore more specific:

\begin{enumerate}[label=(\roman*)]
    \item
    to emphasize and exploit the fact that \emph{arbitrary}
    simultaneous order-statistic CDF bands induce affine constraints in
    two-parameter location--scale families;

    \item
    to give an exact computational treatment of the resulting
    polyhedral confidence region, including marginal projections and
    area;

    \item
    to exploit monotonicity of standard confidence bands for efficient
    envelope construction;

    \item
    to compare different distribution-free CDF bands according to
    induced \emph{parametric} criteria rather than CDF-band width or
    goodness-of-fit power;

    \item
    to investigate systematically how the ranking of confidence bands
    depends on both the inferential target and the base
    location--scale family.
\end{enumerate}

A complete literature review is still required before any stronger
novelty claim is made, especially for the general polyhedral and
algorithmic statements.


\section{A parameter-targeted band optimization problem}

The simulation results suggest a broader optimization problem.

Let $\cB_\alpha$ denote a class of simultaneous distribution-free CDF
bands with coverage at least $1-\alpha$.  Every
$B\in\cB_\alpha$ induces a confidence region
\[
    \cC_B(X)
\]
for $(\mu,\sigma)$.

A band that is optimal as a confidence band for the entire CDF need not
be optimal after transformation into parameter space.  This motivates,
for a fixed location--scale family $G$, the criterion
\begin{equation}
    B_A^\star
    =
    \arg\min_{B\in\cB_\alpha}
    \Ee_{0,1}
    \left[
        \operatorname{Area}
        \left(
            \cC_B(X)
        \right)
    \right].
    \label{eq:optarea}
\end{equation}

Likewise one may define
\begin{align}
    B_\mu^\star
    &=
    \arg\min_{B\in\cB_\alpha}
    \Ee_{0,1}
    \left[
        |\cC_{\mu,B}(X)|
    \right],
    \label{eq:optmu}
    \\
    B_\sigma^\star
    &=
    \arg\min_{B\in\cB_\alpha}
    \Ee_{0,1}
    \left[
        |\cC_{\sigma,B}(X)|
    \right].
    \label{eq:optsigma}
\end{align}

The preliminary simulations provide direct evidence that these three
objectives need not agree:
\[
    B_A^\star
    \neq
    B_\mu^\star
    \neq
    B_\sigma^\star
\]
within the set of bands currently compared.

This suggests that ``optimal CDF band'' is not an intrinsic concept
once the CDF band is being used as an intermediate device for
parametric inference.  Optimality should instead be defined relative to
the parameter target.


\section{Further work}

Several extensions are natural.

\subsection{Larger simulation design}

The first priority is a systematic experiment over
\[
    n\in
    \{20,50,100,200,500,1000\}
\]
and several base distributions, for example
\[
    \text{normal},
    \quad
    \text{logistic},
    \quad
    \text{Laplace},
    \quad
    t_2,t_3,t_5,t_{10},
\]
and potentially asymmetric location--scale families.

The principal plots should display relative criteria
\begin{align}
    \frac{
        \Ee[A_B]
    }{
        \Ee[A_{\mathrm{KS}}]
    },
    \qquad
    \frac{
        \Ee[L_{\mu,B}]
    }{
        \Ee[L_{\mu,\mathrm{KS}}]
    },
    \qquad
    \frac{
        \Ee[L_{\sigma,B}]
    }{
        \Ee[L_{\sigma,\mathrm{KS}}]
    }
\end{align}
as functions of sample size and tail behavior.

Because every method is applied to the same simulated sample, paired
comparisons should also be reported, for example
\[
    \Pp
    \left(
        A_{\mathrm{DW}}
        <
        A_{\mathrm{KS}}
    \right).
\]


\subsection{Coverage matching}

Native calibrations do not always have identical finite-sample
coverage.  DKW is deliberately conservative.  The rank-calibrated KS,
BJ, and DW thresholds have an explicit marginal coverage
$k/(B+1)$, although their conditional coverage given one realized
calibration sample still fluctuates around that level.

Two comparisons are therefore useful:

\begin{enumerate}[label=(\roman*)]
    \item
    the procedures under their native published calibration;

    \item
    a coverage-matched comparison in which thresholds are adjusted to
    give a common finite-sample coverage as close as possible to
    $1-\alpha$.
\end{enumerate}

The second comparison isolates geometric efficiency from differences in
conservativeness.


\subsection{Asymptotic theory}

The striking stability of some normal-theory ratios suggests that
asymptotic relative efficiencies may exist.  For example, the
Berk--Jones scale interval was approximately $38\%$ shorter than the KS
scale interval throughout the preliminary range $n=20$ to $200$.

A theoretical question is whether quantities such as
\[
    \frac{
        \Ee[
            |\cC_{\sigma,\mathrm{BJ}}|
        ]
    }{
        \Ee[
            |\cC_{\sigma,\mathrm{KS}}|
        ]
    }
\]
converge, and how the limit depends on the tail behavior of
$G^{-1}$.


\subsection{Optimization over bands}

The most ambitious extension is to optimize the band limits
$(\ell_i,u_i)$ directly subject to the distribution-free coverage
constraint
\[
    \Pp
    \left(
        \ell_i
        \leq
        U_{(i)}
        \leq
        u_i
        \text{ for every }i
    \right)
    \geq
    1-\alpha.
\]

Frey \cite{Frey2008} optimizes CDF bands according to CDF-scale
narrowness criteria.  The present framework suggests replacing that
objective by expected induced parameter-region area or marginal
length.  The resulting optimal band would depend on the base family
$G$ and on the inferential target.


\subsection{Robustness and misspecification}

A further extension is to distinguish confidence-region efficiency
under the assumed model from behavior under misspecification.  If data
are generated from a distribution outside the chosen location--scale
family, the confidence region may become empty.  The event
\[
    \cC(X)=\varnothing
\]
then has a natural model-compatibility interpretation: no member of the
specified family lies inside the chosen simultaneous CDF band.

Thus the same geometric object can potentially combine parametric
uncertainty and goodness-of-fit information.


\section{Conclusion}

A simultaneous distribution-free CDF band can be viewed not only as a
nonparametric confidence object but also as a mechanism for generating
finite-sample confidence regions within a parametric family.

For a continuous two-parameter location--scale family, this inversion
has a particularly simple structure.  Every order-statistic band
constraint becomes affine in $(\mu,\sigma)$, so the joint confidence
region is convex and polyhedral.  Its scale projection, location
projection, vertices, and area can be computed exactly using
piecewise-linear envelopes rather than parameter grids or repeated
simulation over candidate parameter values.

This geometry provides a common framework in which different CDF bands
can be compared according to the uncertainty they induce about
parameters.  Preliminary results demonstrate that such comparisons are
not trivial.  Under normality, KS is most efficient for location,
Berk--Jones is most efficient for scale, and Dümbgen--Wellner is most
efficient for joint area among the bands considered.  The same broad
pattern appears for Laplace data, while sufficiently heavy tails alter
some rankings.

The emerging conclusion is that the best distribution-free CDF band
for parametric inference depends on what parameter uncertainty is being
optimized and on the geometry of the underlying quantile function.
This motivates the study of parameter-targeted distribution-free bands
as a distinct inferential problem.


\appendix

\section{Algorithmic summary}

Given a sample $x_1,\ldots,x_n$, band bounds
$(\ell_i,u_i)$, and a base quantile function $g=G^{-1}$:

\begin{enumerate}
    \item
    Sort the data:
    \[
        x_{(1)}\leq\cdots\leq x_{(n)}.
    \]

    \item
    Transform the probability limits:
    \[
        a_i=g(\ell_i),
        \qquad
        b_i=g(u_i).
    \]

    \item
    Form the lower envelope
    \[
        L(s)
        =
        \max_i
        \{x_{(i)}-b_i s\}.
    \]

    \item
    Form the upper envelope
    \[
        U(s)
        =
        \min_i
        \{x_{(i)}-a_i s\}.
    \]

    \item
    Merge the two ordered breakpoint lists, evaluate
    \[
        H(s)=U(s)-L(s),
    \]
    and locate its first and last zero crossings.  These are the
    endpoints of the feasible scale interval.

    \item
    Compute
    \[
        \mu_L=\min L(s),
        \qquad
        \mu_U=\max U(s)
    \]
    over the feasible scale interval.

    \item
    Merge the envelope breakpoints and compute
    \[
        \operatorname{Area}(\cC)
        =
        \int
        [U(s)-L(s)]\,ds
    \]
    exactly interval by interval.
\end{enumerate}


\section{Monte Carlo calibration}

For KS, BJ, and DW, the critical values can be calibrated once for each
$(n,\alpha)$ using independent uniform samples.  These critical values
and the resulting probability-scale bounds are then reused in every
outer Monte Carlo replication and for every base location--scale
family.

The reference code implements the exchangeable-rank choice
$k=\lceil(B+1)(1-\alpha)\rceil$ directly.  Each simulation function has
an explicit seed argument, so a displayed result can be reproduced
without adding random-number-state machinery to the main algorithm.

This separation is computationally important:
\[
    \text{band calibration}
    \quad\longrightarrow\quad
    (\ell_i,u_i)
    \quad\longrightarrow\quad
    \text{family-specific quantile transformation}.
\]

Changing from the normal family to Laplace or Student-$t$, for example,
requires changing only
\[
    G^{-1},
\]
not recalibrating the distribution-free band.


\section{Reference R implementation}
\label{app:Rcode}

The following commented R implementation contains the complete
computational pipeline used in this note: Monte Carlo calibration of the
KS, Berk--Jones, and D\"umbgen--Wellner bands; construction of the DKW
bounds; exact envelope-based computation of the induced location--scale
confidence region; and the outer Monte Carlo comparison.  The final
driver accepts an arbitrary base quantile function, so the same code
applies to the normal, Laplace, and Student-$t$ experiments.

The listing is intentionally pedagogical rather than package-level
production code.  It assumes that arguments have the advertised types
and ranges and omits routine checks for malformed inputs.  It does retain
the mathematical edge cases that affect the result---empty and
unbounded confidence sets, infinite transformed probability endpoints,
and relative numerical tolerances after affine normalization.  This
keeps the statistical construction visible while the companion
verification script tests correctness independently.

The listing is consolidated from the accompanying notebook
\path{location_scale_research.ipynb}.  It retains the explanatory
comments from the notebook.  Sourcing the file only defines functions:
the illustrative example and the complete reproduction study are
explicit drivers and are not launched automatically.  The latter
calibrates each required sample size once and reuses those bands across
the normal, Laplace, and Student-$t_2$ experiments.  Notebook outputs are
omitted; the numerical summaries reported in the main text remain the
record of the preliminary runs.

\lstinputlisting[
    style=Rcode,
    caption={Pedagogical commented R implementation and reproduction driver.},
    label={lst:reference-R}
]{location_scale_reference.R}


\section{Verification code}
\label{app:Rchecks}

The separate companion file \path{location_scale_checks.R} checks the
implementation against the independent $O(n^2)$ pairwise formula on
randomized inputs, direct envelope maxima and minima, and dense
numerical area integration.  It also contains regression tests for
one-segment and final-segment envelope evaluation; extreme affine
changes of units; empty and unbounded confidence sets; equivalence
between each inverted band and its originating statistic event; rank
calibration; seeded reproducibility; and reuse of precomputed bands.
Finally, an empirical log--log timing slope guards against an accidental
return to quadratic geometry.  The test code is distributed with the
notebook but is not reproduced here, so the printed appendix remains
focused on the algorithm a reader is meant to study.


\begin{thebibliography}{99}
\small

\bibitem{BarrDavidson1973}
D. R. Barr and T. Davidson.
A Kolmogorov--Smirnov test for censored samples.
\emph{Technometrics},
15(4):739--757, 1973.
doi:10.1080/00401706.1973.10489108.

\bibitem{BerkJones1979}
R. H. Berk and D. H. Jones.
Goodness-of-fit test statistics that dominate the Kolmogorov statistics.
\emph{Zeitschrift f\"ur Wahrscheinlichkeitstheorie und Verwandte Gebiete},
47(1):47--59, 1979.
doi:10.1007/BF00533250.

\bibitem{Chattopadhyay2026}
S. Chattopadhyay, S. Sarkar, and A. K. Kuchibhotla.
Inference for quantile-parametrized families via CDF confidence bands.
In \emph{Proceedings of the 42nd Conference on Uncertainty in
Artificial Intelligence},
Proceedings of Machine Learning Research, vol. 337,
pp. 1034--1053, 2026.

\bibitem{DKW1956}
A. Dvoretzky, J. Kiefer, and J. Wolfowitz.
Asymptotic minimax character of the sample distribution function and
of the classical multinomial estimator.
\emph{Annals of Mathematical Statistics},
27(3):642--669, 1956.

\bibitem{DumbgenWellner2023}
L. D\"umbgen and J. A. Wellner.
A new approach to tests and confidence bands for distribution
functions.
\emph{Annals of Statistics},
51(1):260--289, 2023.
doi:10.1214/22-AOS2249.

\bibitem{Easterling1976}
R. G. Easterling.
Goodness of fit and parameter estimation.
\emph{Technometrics},
18(1):1--9, 1976.
doi:10.1080/00401706.1976.10489394.

\bibitem{Frey2008}
J. C. Frey.
Optimal distribution-free confidence bands for a distribution
function.
\emph{Journal of Statistical Planning and Inference},
138(10):3086--3098, 2008.
doi:10.1016/j.jspi.2007.12.001.

\bibitem{FreyMarreroNorton2009}
J. C. Frey, O. Marrero, and D. E. Norton.
Minimum-area confidence sets for a normal distribution.
\emph{Journal of Statistical Planning and Inference},
139(3):1023--1032, 2009.
doi:10.1016/j.jspi.2008.06.005.

\bibitem{LittellRao1978}
R. C. Littell and P. V. Rao.
Confidence regions for location and scale parameters based on the
Kolmogorov--Smirnov goodness of fit statistic.
\emph{Technometrics},
20(1):23--27, 1978.

\bibitem{Massart1990}
P. Massart.
The tight constant in the Dvoretzky--Kiefer--Wolfowitz inequality.
\emph{Annals of Probability},
18(3):1269--1283, 1990.
doi:10.1214/aop/1176990746.

\bibitem{Owen1995}
A. B. Owen.
Nonparametric likelihood confidence bands for a distribution function.
\emph{Journal of the American Statistical Association},
90(430):516--521, 1995.
doi:10.1080/01621459.1995.10476543.

\bibitem{Salvia1980}
A. A. Salvia.
Some fundamental properties of Kolmogorov--Smirnov consonance sets.
\emph{Technometrics},
22(1):109--111, 1980.
doi:10.1080/00401706.1980.10486107.

\end{thebibliography}

\end{document}
